%!TEX root = ../thesis.tex
%*******************************************************************************
%****************************** Discussion *************************************
%*******************************************************************************

\chapter{Discussion}
\label{ch:discussion}

%********************************** Section 4.1 ********************************
\section{Comparison with Existing Systems}

The developed platform successfully addresses the limitations of the first-generation ``Artery in a Box'' prototype \citep{Kacerik2024}, which relied on manual adjustment of pump parameters without feedback control. The original system required an operator to monitor flow rate readings and manually tune the drive voltage, a process that was both labour-intensive and unable to compensate for the slow thermal and hydraulic drifts demonstrated in Section~\ref{sec:open_vs_closed}. By introducing automated PID regulation, the present platform eliminates operator dependence and ensures consistent shear conditions over extended experimental durations.

Compared to commercial syringe pump systems commonly used in microfluidic research \citep{Whitesides2006, Squires2005}, the micropump-based platform offers several distinct advantages. The total component cost is estimated under \pounds200, compared to \pounds1000--5000 for commercial syringe pumps with equivalent flow control capabilities. The compact form factor (the entire fluidic assembly fits within a \SI{20}{\centi\metre} $\times$ \SI{15}{\centi\metre} footprint) enables integration into standard incubator environments for live-cell experiments. Most importantly, the recirculating architecture allows NPs to experience repeated passes through the test section, mimicking the circulation of blood through the vasculature --- a feature fundamentally unavailable in single-pass syringe pump configurations.

However, syringe pumps offer advantages in applications requiring precise volumetric delivery or operation with viscous fluids. The piezoelectric micropump's flow-pressure characteristic is inherently more sensitive to downstream resistance changes than a positive-displacement syringe, making closed-loop control essential rather than optional for this platform.

%********************************** Section 4.2 ********************************
\section{Significance for Nanocarrier Research}

The ability to maintain stable, physiologically relevant shear conditions is critical for meaningful NP stability studies. Previous work has demonstrated that NP behaviour under static conditions poorly predicts performance in flowing blood \citep{Cooley2018, Thomas2014}, as shear forces influence protein corona formation, NP aggregation, and cellular uptake mechanisms \citep{Blanco2015}. The shear rate mapping presented in Table~\ref{tab:shear_mapping} confirms that the platform can replicate conditions spanning from venous ($\sim$\SI{30}{\per\second}) to arterial ($\sim$\SI{1300}{\per\second}) flow regimes, enabling systematic investigation of shear-dependent NP behaviour.

The closed-loop stability demonstrated in Section~\ref{sec:open_vs_closed} is particularly significant. NP stability studies typically require continuous perfusion over \SIrange{15}{60}{\minute} to observe meaningful changes in particle size distribution or surface properties \citep{Mitchell2021}. The platform's ability to maintain flow within $\pm$2\% of the target ensures that observed changes in NP characteristics can be attributed to biological or physicochemical processes rather than experimental artefacts from flow instability.

The real-time shear rate computation displayed in the GUI represents a practical advantage over systems where researchers must perform post-hoc calculations. By directly monitoring the haemodynamic conditions during experiments, operators can verify that physiological relevance is maintained throughout the test and respond immediately to any deviations.

%********************************** Section 4.3 ********************************
\section{Limitations}

Several limitations of the current system should be acknowledged. First, all characterisation was performed using deionised water. Biological media containing proteins and NPs will have different viscosity and surface tension properties that may alter the pump's flow-pressure characteristics and the sensor's calibration accuracy. The SLF3S is factory-calibrated for water and IPA only; fluids with significantly different thermal properties may require custom calibration coefficients. Re-tuning of PID gains may also be necessary when switching to higher-viscosity blood-analogue fluids.

Second, the \SI{10}{\hertz} control bandwidth, while adequate for the slow dynamics of the current fluidic circuit (where the dominant time constants are on the order of seconds), may be insufficient for systems with faster pressure transients, such as those incorporating active microvalves or rapidly switching flow paths.

Third, the current single-channel design limits experimental throughput. Many organ-on-chip applications require parallel flow channels with independent control to enable simultaneous testing of multiple conditions \citep{Bhatia2014, Sosa-Hernandez2018}. The ESP32's multiple I\textsuperscript{2}C buses and PWM channels could support multi-channel operation in principle, but the current firmware and GUI are designed for single-channel use only.

Fourth, the system has not yet been validated with cultured cells or in the presence of NPs. The biocompatibility of the fluidic pathway (PPSU pump body, silicone tubing, glass sensor) is expected to be adequate based on material specifications, but this requires experimental confirmation.

%********************************** Section 4.4 ********************************
\section{Future Work}

Future development will focus on three directions. The primary objective is integrating the platform with microfluidic chips containing cultured endothelial cells to study NP adhesion and uptake under controlled shear conditions --- the intended application that motivated this project. This will require validation of cell viability under recirculation conditions and development of protocols for introducing and recovering NP samples.

Second, replacing deionised water with viscosity-matched blood-analogue fluids (such as glycerol-water mixtures or xanthan gum solutions) will better replicate whole blood rheology \citep{Popel2005}. This will necessitate recalibration of the flow sensor and re-tuning of PID parameters, but the platform's architecture --- with software-adjustable gains and real-time monitoring --- is well suited to accommodate such changes.

Third, extending the platform to multi-channel operation would significantly increase experimental throughput. This could be achieved by multiplexing pump and sensor arrays on a single ESP32, leveraging its multiple I\textsuperscript{2}C buses and sixteen PWM channels, combined with I\textsuperscript{2}C address multiplexers for sensor arrays sharing the same default address.
